\section{Conclusion}
\label{sec:conclusion}
In this paper, we present the first comprehensive survey of verbal reinforcement learning, organized around three pillars based on when language enters an agent's lifecycle: grounding signal, deliberative feedback, and learning signal. Together, these pillars reveal a common pattern: verbal feedback is rapidly becoming the primary medium for defining, updating, and improving agents.

Looking ahead, we anticipate that verbal feedback will become a first-class
design primitive in agent architectures. The boundaries between pillars will
increasingly blur as agents consume verbal feedback in unified loops that span
grounding, deliberation, and learning within a single trajectory. The bottleneck
will shift from \textit{generating} feedback to \textit{verifying} it, making
feedback provenance, quality benchmarks, and adversarial robustness first-class
infrastructure requirements. Finally, tool interfaces will need to be redesigned
around agent consumability rather than human readability, as the gap between
tool output quality and agent performance becomes the practical ceiling for
iterative refinement.


% In this survey, we present the first comprehensive study of verbal reinforcement learning, organizing it around three pillars distinguished by when language enters an agent's lifecycle: grounding signal, deliberative feedback, and learning signal. Verbal RL has already driven improvements across coding assistants, scientific discovery, robotics, and other domains where these pillars increasingly operate in concert; we discuss applications in detail in Appendix~\ref{app:applications}. As LLMs continue to improve at generating and interpreting feedback, verbal feedback will expand as the primary medium through which agents are corrected, aligned, and grounded. However, effectiveness depends critically on feedback quality: progress requires feedback-specialized models, benchmarks for evaluating feedback utility, and agent-centered tool interfaces.

% Looking ahead, we anticipate that verbal feedback will become a first-class design primitive in agent architectures, with dedicated feedback models, standardized quality metrics, and agent-centric tool APIs emerging as core infrastructure for building more capable, adaptable, and aligned agents.

% Progress requires feedback-specialized models, benchmarks for evaluating feedback utility, and agent-centered tool interfaces whose outputs are easier to interpret and act upon. In particular, the field needs stronger methods for assessing whether feedback is precise, actionable, and robust, alongside compatibility metrics that measure how effectively agents consume tool outputs.

 
% \das{do we want to end with some type of a forecast?}


 % However, effectiveness depends critically on feedback quality: progress requires feedback-specialized models, benchmarks for evaluating feedback utility, and agent-centered tool interfaces. We anticipate that verbal feedback will become a first-class design primitive, with dedicated feedback models, standardized quality metrics, and agent-centric tool APIs emerging as core infrastructure.